\documentclass[11pt,a4paper]{article}
\usepackage[utf8]{inputenc}
\usepackage[T1]{fontenc}
\usepackage{lmodern}
\usepackage[margin=1.5cm]{geometry}
\usepackage{multicol}
\usepackage{enumitem}
\usepackage{microtype}
\usepackage{xcolor}
\usepackage{titlesec}
\usepackage[hidelinks]{hyperref}

\definecolor{heading}{HTML}{164B70}
\titleformat{\section}{\normalsize\bfseries\color{heading}}
  {\thesection.}{0.4em}{}
\titlespacing*{\section}{0pt}{7pt}{3pt}
\setlength{\parindent}{0pt}
\setlength{\parskip}{3pt}
\setlength{\columnsep}{7mm}
\setlength{\multicolsep}{4pt}
\setlist[itemize]{leftmargin=*,nosep,topsep=2pt}
\setlist[enumerate]{leftmargin=*,nosep,topsep=2pt}
\setlength{\emergencystretch}{1.5em}
\urlstyle{same}
\pagestyle{plain}

\begin{document}

\begin{center}
{\Large\bfseries Cybersecurity and Digital Forensics}\\[2pt]
%{\small Two-page revision notes}
{\large\bfseries Harrison Katantazi}
\end{center}
\vspace{-6pt}

% ==================== PAGE 1 ====================
\begin{multicols}{2}
\raggedcolumns

\section{Cyber Threat Sources and Actors}

A \textbf{threat source} is a person, condition, or circumstance
that could harm information systems. NIST identifies four
categories [1]:
\begin{itemize}
\item \textbf{Adversarial:} Deliberate harmful activity.
Example: criminals steal customer records.
\item \textbf{Accidental:} Unintentional human errors.
Example: an administrator exposes confidential cloud files.
\item \textbf{Structural:} Hardware, software, or equipment
failures. Example: a failed disk makes data unavailable.
\item \textbf{Environmental:} Natural disasters or external
infrastructure failures. Example: flooding or a power outage.
\end{itemize}

A \textbf{threat actor} is an individual or group deliberately
conducting harmful cyber activity. Actors are one type of
threat source. Common categories include [2]:
\begin{itemize}
\item \textbf{Cybercriminals:} Seek money; for example,
a ransomware group demanding payment.
\item \textbf{State-sponsored actors:} Pursue government
objectives; for example, stealing sensitive research.
\item \textbf{Hacktivists:} Pursue political or ideological
goals; for example, disrupting a website in protest.
\item \textbf{Malicious insiders:} Misuse authorised access;
for example, an employee selling confidential records.
\item \textbf{Thrill-seekers:} Seek recognition or satisfaction;
for example, breaking into a website to demonstrate ability.
\item \textbf{Terrorist groups:} Seek intimidation or ideological
violence; for example, attempting to disrupt critical services.
\end{itemize}

\textbf{Distinction:} Accidental database deletion is an
accidental threat source; deliberate deletion for revenge
involves a malicious insider threat actor.

\section{Code of Practice for Information Security}

\textbf{ISO/IEC 17799} specified the \textit{Code of Practice
for Information Security Management}. It was renumbered
\textbf{ISO/IEC 27002}. The current edition is
\textbf{ISO/IEC 27002:2022}, \textit{Information Security,
Cybersecurity and Privacy Protection---Information Security
Controls}. It provides control guidance; ISO/IEC 27001
specifies ISMS requirements [3].

\columnbreak

\section{Chain of Custody: RFC 3227}

The \textbf{chain of custody} is the documented history of
evidence: how it was discovered, handled, stored, and
transferred. Section~4.1 specifies recording [4]:
\begin{enumerate}
\item \textbf{Discovery and collection:} Where, when,
and by whom evidence was found and collected.
\item \textbf{Handling and examination:} Where, when,
and by whom it was handled or examined.
\item \textbf{Custody and storage:} Each custodian,
the period of custody, and storage arrangements.
\item \textbf{Transfers:} When and how custody changed,
including shipping or tracking numbers where applicable.
\end{enumerate}

Evidence must be securely protected. Section~4.2 calls for
strictly restricted, documented access and the ability to
detect unauthorised access. These measures support
traceability, authenticity, and reliable handling.

\textbf{Example:} Record who collected a drive, who examined
it, where it was secured, and every handover.

\section{Forensics in Incident Response}

Forensic techniques preserve evidence, establish what
happened, and guide response decisions [5].
\begin{enumerate}
\item \textbf{Prepare:} Train responders, configure logging,
prepare tools, and define evidence handling procedures.
\item \textbf{Collect and preserve:} Acquire relevant logs,
network traffic, memory, and disk images. Prioritise volatile
data where appropriate, verify integrity using hashes,
and maintain a chain of custody.
\item \textbf{Examine and analyse:} Extract relevant information,
reconstruct the timeline, investigate the entry point,
and assess affected systems and data.
\item \textbf{Contain and recover:} Use findings to isolate
compromised systems, remove malicious components, address
vulnerabilities, and support safe restoration.
\item \textbf{Report and improve:} Document methods,
findings, limitations, and actions; improve controls and
response procedures.
\end{enumerate}

\textbf{Operational balance:} Preserve evidence while
considering the urgency of stopping harm and maintaining
essential services.

\textbf{Example:} During ransomware response, preserve
available memory and logs, analyse forensic copies,
identify affected systems, and guide containment and recovery.

\end{multicols}

\newpage

% ==================== PAGE 2 ====================
\begin{multicols}{2}
\raggedcolumns

\section{NIS 2 Directive}

\textbf{Directive (EU) 2022/2555} strengthens cybersecurity
across the EU and replaces NIS~1. NIS refers to
\textit{Network and Information Systems} [6].
\begin{itemize}
\item \textbf{Scope:} Essential and important entities in
sectors including energy, transport, healthcare, water,
and digital infrastructure; generally medium-sized and
large entities, with exceptions.
\item \textbf{Risk management:} Incident handling, business
continuity, supply chain security, and vulnerability management.
\item \textbf{Accountability:} Management approves and
oversees cybersecurity risk-management measures.
\item \textbf{Reporting:} Significant incidents generally
require an early warning within \textbf{24 hours} of awareness,
notification within \textbf{72 hours}, and a final report
within \textbf{one month after notification}. Special rules
apply to trust service providers and ongoing incidents.
\item \textbf{Implementation:} Member States transpose
the directive into national law, with supervision,
enforcement, and cross-border cooperation.
\end{itemize}
\textbf{Example:} A covered hospital must manage cyber risks
and report significant incidents.

\section{Cyber Resilience Act}

\textbf{Regulation (EU) 2024/2847} establishes cybersecurity
requirements for products with digital elements supplied
commercially on the EU market. It generally covers hardware
and software connecting directly or indirectly to devices
or networks, subject to exclusions [7].
\begin{itemize}
\item \textbf{Security by design:} Assess product risks
and integrate appropriate protections.
\item \textbf{Ongoing support:} Handle vulnerabilities,
provide security updates, and state the support end date.
\item \textbf{Compliance:} Technical documentation,
conformity assessment, and CE marking; some categories
require independent assessment.
\item \textbf{Reporting:} Manufacturers report actively
exploited vulnerabilities and severe product-security
incidents: early warning within \textbf{24 hours} and
notification within \textbf{72 hours} of awareness.
\end{itemize}
\textbf{Dates:} Entered into force 10 December 2024;
manufacturer reporting began 11 September 2026;
main requirements apply 11 December 2027.

\textbf{Example:} A connected-camera manufacturer must
address security risks and provide support-period updates.

\columnbreak

\section{Zero Trust: NIST SP 800-207}

Zero Trust removes automatic trust based on network
location or ownership. Access depends on verified identity,
device security, and policy [8].

\textbf{Seven principles:}
\begin{enumerate}
\item Treat data sources and computing services as resources.
\item Secure communications on every network.
\item Authorise each session with least privilege.
\item Use dynamic, context-aware access policies.
\item Monitor asset integrity and security posture.
\item Authenticate and authorise before access;
reevaluate ongoing sessions according to policy.
\item Use asset, network, and access information to
improve security.
\end{enumerate}

\textbf{Components:} The \textbf{Policy Engine} decides
whether access is allowed; the \textbf{Policy Administrator}
executes decisions; the \textbf{Policy Enforcement Point}
enables, monitors, and terminates connections.

\textbf{Example:} An office employee still needs authorisation
and a sufficiently secure device to access a sensitive database.

\section{Adversarial ML: NIST AI 100-2e2025}

This March~2025 report provides shared terminology for
attacks and mitigations affecting predictive and generative
AI. It classifies threats by lifecycle stage, attacker
objectives, capabilities, knowledge, and data modality [9].

\textbf{Predictive AI attacks:}
\begin{itemize}
\item \textbf{Evasion:} Manipulated inputs cause incorrect
predictions during use.
\item \textbf{Poisoning:} Corrupted training data or models
alter behaviour, potentially introducing backdoors.
\item \textbf{Privacy attacks:} Extract sensitive data or
information about the model.
\end{itemize}

\textbf{Generative AI:} Covers supply chain poisoning,
direct prompting attacks, indirect prompt injection through
external content, and AI agent security.

\textbf{Mitigations:} Adversarial training, data sanitisation,
model inspection, differential privacy, and prompt-injection
detection.

\textbf{Limitations:} Defences involve trade-offs in accuracy,
privacy, robustness, and cost, and require rigorous evaluation
against capable attackers. The report offers voluntary
guidance for understanding and assessing AI security risks.

\end{multicols}

\vspace{1pt}
\hrule
\vspace{3pt}
{\footnotesize
\textbf{Official references} \quad
\href{https://doi.org/10.6028/NIST.SP.800-30r1}{[1] NIST SP 800-30 Rev.~1};
\href{https://www.cyber.gc.ca/en/guidance/introduction-cyber-threat-environment}{[2] Canadian Centre for Cyber Security: Threat Environment};
\href{https://www.iso.org/standard/50297.html}{[3] ISO/IEC 27002 (history)}
and \href{https://www.iso.org/standard/75652.html}{2022 edition};
\href{https://www.rfc-editor.org/rfc/rfc3227.html}{[4] RFC 3227, Sections 4.1--4.2};
\href{https://doi.org/10.6028/NIST.SP.800-86}{[5] NIST SP 800-86};
\href{https://digital-strategy.ec.europa.eu/en/policies/nis2-directive}{[6] European Commission: NIS 2};
\href{https://digital-strategy.ec.europa.eu/en/policies/cra-summary}{[7] European Commission: CRA Summary}
and \href{https://digital-strategy.ec.europa.eu/en/policies/cra-reporting}{Reporting};
\href{https://doi.org/10.6028/NIST.SP.800-207}{[8] NIST SP 800-207};
\href{https://doi.org/10.6028/NIST.AI.100-2e2025}{[9] NIST AI 100-2e2025}.
}

\end{document}